% Part: methods
% Chapter: induction
% Section: inductive-definitions

\documentclass[../../../include/open-logic-section]{subfiles}

\begin{document}

\olfileid{mth}{ind}{idf}

\olsection{استقرائی تعریفاں}

منطق وچ اسیں اکثر شےواں دیاں قسماں \emph{استقرائی طریقے نال} تعریف کردے آں۔ یعنی اسیں قاعدے دیندے آں کہ کیہڑی شے ایس قسم دی منّی جائے، تے قاعدے سمجھاؤندے نیں کہ اوس قسم دیاں پرانیاں شےواں توں اوسے قسم دیاں نویاں شےواں کیویں بنن۔ مثلاً اسیں نشاناں دیاں خاص لڑیواں، جیویں کسے زبان دے ٹرم تے !!{formula}s، اکثر استقرا نال تعریف کردے آں۔ اک سوکھی مثال لئی اوہ لڑیواں ویکھو جیہڑیاں اکھراں $\mathrm{a}$، $\mathrm{b}$، $\mathrm{c}$، $\mathrm{d}$، نشان~$\circ$، تے بریکٹاں $[$ تے $]$ نال بندیاں نیں، جیویں «$[[\mathrm{c} \circ \mathrm{d}][$»، «$[\mathrm{a}[]\circ]$»، «$\mathrm{a}$»، یا «$[[\mathrm{a} \circ \mathrm{b}]\circ
\mathrm{d}]$»۔ تہانوں شاید لگے کہ پہلیاں دوہاں لڑیواں وچ کوئی «غلطی» اے۔ پہلی وچ بریکٹاں دا میل بالکل نہیں، تے تہانوں لگ سکدا اے کہ «$\circ$» نوں اوہ عبارتاں «جوڑنیاں» چاہیدیاں نیں جیہڑیاں آپ ٹھیک ہون۔ تیجی تے چوتھی لڑی ٹھیک لگدیاں نیں۔ ہر «$[$» لئی اک بند کرن والا «$]$» اے (جے کوئی بریکٹ ہے وی)، تے ہر $\circ$ دے دوہاں پاسیاں اُتے «ٹھیک» عبارتاں لبھدیاں نیں، جیہناں دے دوالے بریکٹاں دی اک جوڑی اے۔

اسیں این ٹھیک دسنا چاہندے آں کہ «ٹھیک ٹرم» کیہنوں آکھیے۔ پہلاں، ہر اکھر آپے ٹھیک اے۔ جیہڑی شے بس اکلّا اکھر نہیں، اوہ «$[t \circ s]$» دی صورت وچ ہووے، جتھے $s$ تے $t$ آپ ٹھیک ہون۔ اُلٹے پاسے، جے $t$ تے $s$ ٹھیک نیں، تاں اوہناں دے وچکار $\circ$ رکھ کے تے دوالے بریکٹاں دی اک جوڑی لا کے اسیں نواں ٹھیک ٹرم بنا سکدے آں۔ ایہ عمل ورت کے اسیں ٹھیک ٹرماں دا سیٹ \emph{تعریف} کر سکدے آں۔ ایہ \emph{استقرائی تعریف} اے۔

\begin{defn}[ٹھیک ٹرم]
  \emph{ٹھیک ٹرماں} دا سیٹ استقرائی طریقے نال ایویں تعریف ہُندا اے:
  \begin{enumerate}
  \item کوئی وی اکھر $\mathrm{a}$، $\mathrm{b}$، $\mathrm{c}$، $\mathrm{d}$ اک ٹھیک ٹرم اے۔
  \item جے $s_1$ تے $s_2$ ٹھیک ٹرم نیں، تاں $[s_1 \circ s_2]$ وی ٹھیک ٹرم اے۔
  \item ہور کوئی شے ٹھیک ٹرم نہیں۔
  \end{enumerate}
\end{defn}

ایہ تعریف دسدی اے کہ کوئی شے اُدوں تے صرف اُدوں ٹھیک ٹرم اے جدوں اوہ دو شرطاں (۱) تے~(۲) مطابق محدود گنتی دے قدماں وچ بن سکے۔ پہلے قدم وچ اسیں اوہ سارے ٹھیک ٹرم بناؤندے آں جیہڑے بس اکلّے اکھر نیں، یعنی:
\[
\mathrm{a}, \mathrm{b}, \mathrm{c}, \mathrm{d}
\]
دوجے قدم وچ اسیں اپنے بنائے ٹرماں اُتے (۲) لاؤندے آں۔ سانوں ایہ ملن گے:
\[
[\mathrm{a} \circ \mathrm{a}], [\mathrm{a} \circ \mathrm{b}],
[\mathrm{b} \circ \mathrm{a}], \dots, [\mathrm{d} \circ \mathrm{d}]
\]
دو اکھراں دے سارے جوڑاں لئی۔ تیجے قدم وچ اسیں (۲) فیر ہُن تک بنائے کسے وی دو ٹھیک ٹرماں اُتے لاؤندے آں۔ سانوں نویں ٹھیک ٹرم ملدے نیں، جیویں $[\mathrm{a} \circ [\mathrm{a} \circ
    \mathrm{a}]]$، جتھے $t$ پہلے قدم والا $\mathrm{a}$ اے تے $s$ دوجے قدم والا $[\mathrm{a} \circ \mathrm{a}]$ اے، تے $[[\mathrm{b} \circ
    \mathrm{c}] \circ [\mathrm{d} \circ \mathrm{b}]]$، جیہڑا دوجے قدم دے دو ٹرماں $[\mathrm{b} \circ \mathrm{c}]$ تے $[\mathrm{d}
  \circ \mathrm{b}]$ توں بنیا اے۔ تے ایس طرح اگے۔ شق (۳) ایہ روکدی اے کہ جیہڑی شے ایویں نہیں بنی، اوہ ٹھیک ٹرماں دے سیٹ وچ آ وڑے۔

دھیان دیو کہ اسیں ہن تک ایہ ثابت نہیں کیتا کہ نشاناں دی ہر «ٹھیک لگدی» لڑی ایس تعریف مطابق ٹھیک اے۔ پر ایہ صاف ہونا چاہیدا اے کہ اسیں جیہڑی وی لڑی بنا سکدے آں اوہ واقعی «ٹھیک لگدی» اے۔ بریکٹاں دا میل اے تے $\circ$ اوہ حصے جوڑدا اے جیہڑے آپ ٹھیک نیں۔

استقرائی تعریفاں دی خاص گل ایہ اے کہ جے تسی سارے ٹھیک ٹرماں بارے کوئی گل ثابت کرنا چاہندے او، تاں تعریف دسدی اے کہ کیہڑیاں صورتاں ویکھنیاں پَین گئیاں۔ مثلاً جے تہانوں دسیا جائے کہ $t$ اک ٹھیک ٹرم اے، تاں استقرائی تعریف دسدی اے کہ $t$ کیہو جیہا ہو سکدا اے۔ $t$ اک اکھر ہو سکدا اے، یا $[s_1 \circ s_2]$، کسے دو ٹھیک ٹرماں $s_1$ تے~$s_2$ لئی۔ شق (۳) کرکے بس ایہو صورتاں ممکن نیں۔

جدوں استقرائی طریقے نال تعریف ہوئے سیٹ دے سارے رکناں بارے دعویٰ ثابت کرنا ہووے، تاں مضبوط استقرا خاص ضروری ہو جاندا اے۔ مثلاً منّ لو اسیں ثابت کرنا چاہندے آں کہ لمبائی~$n$ والے ہر ٹھیک ٹرم وچ $[$ دی گنتی~$< n/2$ اے۔ ایس نوں سارے~$n$ بارے دعویٰ سمجھ سکدے آں۔ ہر $n$ لئی، $[$ دی گنتی، لمبائی~$n$ والے کسے وی ٹھیک ٹرم وچ، $< n/2$ اے۔

\begin{prop}
  کسے وی $n$ لئی، $[$ دی گنتی، لمبائی~$n$ والے کسے ٹھیک ٹرم وچ، $< n/2$ اے۔
\end{prop}

\begin{proof}
ایہ نتیجہ (مضبوط) استقرا نال ثابت کرن لئی سانوں وکھاؤنا پئے گا کہ ہِیٹھلا شرطیہ دعویٰ سچ اے:
\begin{quote}
  جے ہر $l < k$ لئی، لمبائی~$l$ والے کسے وی ٹھیک ٹرم وچ $< l/2$ وار $[$ آؤندا اے، تاں لمبائی~$k$ والے کسے وی ٹھیک ٹرم وچ $< k/2$ وار $[$ آؤندا اے۔
\end{quote}
ایہ شرطیہ وکھاؤن لئی اوہدا مقدم سچ منّو، یعنی منّ لو کہ کسے وی $l<k$ لئی، لمبائی~$l$ والے ٹھیک ٹرماں وچ $< l/2$ وار $[$ آؤندا اے۔ ایس منّی ہوئی گل نوں استقرائی مفروضہ آکھدے آں۔ اسیں وکھاؤنا چاہندے آں کہ ایہو گل لمبائی~$k$ والے ٹھیک ٹرماں لئی وی سچ اے۔

سو منّ لو $t$ لمبائی~$k$ والا اک ٹھیک ٹرم اے۔ کیوں جے ٹھیک ٹرم استقرائی طریقے نال تعریف ہوئے نیں، ساڈے کول دو صورتاں نیں۔ (۱)~$t$ بس اک اکھر اے، یا (۲)~$t$، $[s_1 \circ s_2]$ اے، کسے ٹھیک ٹرماں $s_1$ تے~$s_2$ لئی۔
\begin{enumerate}
\item $t$ اک اکھر اے۔ تاں $k = 1$ تے $[$ دی گنتی $t$ وچ~$0$ اے۔ کیوں جے $0 < 1/2$، دعویٰ پورا ہُندا اے۔
\item $t$، $[s_1 \circ s_2]$ اے، کسے ٹھیک ٹرماں $s_1$ تے~$s_2$ لئی۔ $l_1$،~$s_1$ دی لمبائی تے $l_2$،~$s_2$ دی لمبائی ہووے۔ تاں لمبائی~$k$، جیہڑی $t$ دی اے، $l_1+l_2+3$ اے ($s_1$ تے $s_2$ دیاں لمبائیاں تے تِن نشان $[$، $\circ$، $]$)۔ کیوں جے $l_1+l_2+3$ ہمیشہ $l_1$ توں وڈا اے، $l_1 < k$۔ ایویں ای $l_2
  < k$۔ ایہدا مطلب اے کہ استقرائی مفروضہ ٹرماں $s_1$ تے~$s_2$ اُتے لاگو ہُندا اے۔ گنتی~$m_1$ جیہڑی $[$ دی $s_1$ وچ اے، $< l_1/2$ اے، تے گنتی~$m_2$ جیہڑی $[$ دی~$s_2$ وچ اے، $< l_2/2$ اے۔

  $[$ دی گنتی $t$ وچ، $[$ دی~$s_1$ والی گنتی تے $[$ دی~$s_2$ والی گنتی دا جوڑ، نال~$1$ اے، یعنی $m_1 + m_2 + 1$۔ کیوں جے $m_1
  < l_1/2$ تے $m_2 < l_2/2$، سانوں ملدا اے:
  \[
  m_1 + m_2 + 1 < \frac{l_1}{2} + \frac{l_2}{2} + 1 = \frac{l_1+l_2+2}{2} < \frac{l_1+l_2+3}{2} = k/2.
  \]
\end{enumerate}
ہر صورت وچ اسیں وکھا دتا اے کہ $[$ دی گنتی $t$ وچ $< k/2$ اے (استقرائی مفروضے دی بنیاد اُتے)۔ مضبوط استقرا نال قضیہ نکلدا اے۔
\end{proof}

\begin{prob}
  خاص ٹھیک ٹرماں دا سیٹ ایویں تعریف کرو:
  \begin{enumerate}
  \item کوئی وی اکھر $\mathrm{a}$، $\mathrm{b}$، $\mathrm{c}$، $\mathrm{d}$ اک خاص ٹھیک ٹرم اے۔
  \item جے $s$ اک خاص ٹھیک ٹرم اے، تاں $[s]$ وی خاص ٹھیک ٹرم اے۔
  \item جے $s_1$ تے $s_2$ خاص ٹھیک ٹرم نیں، تاں $[s_1 \circ s_2]$ اک خاص ٹھیک ٹرم اے۔
  \item ہور کوئی شے خاص ٹھیک ٹرم نہیں۔
  \end{enumerate}
  وکھاؤ کہ $[$ دی گنتی، کسے خاص ٹھیک ٹرم~$t$ وچ جیہدی لمبائی~$n$ اے، $\le n/2 +1$ اے۔
\end{prob}

\end{document}
